\section{Introduction}
Population dynamics play a crucial role in shaping the economic, social, and developmental trajectory of any nation. By studying a country's growth trends, dependency ratios, and demographic patterns, we can gain a deeper understanding of the opportunities and challenges it faces. This project focuses on analyzing the population trends of Bangladesh, evaluating its year-on-year growth, demographic dividend, and dependency ratios in comparison with global averages. Furthermore, the study examines how historical data on resources and crop yields can be linked with classical population theories such as Malthusian, Boserup, and the Demographic Transition Theory (DTT). Through this, we aim to assess the country’s stage in demographic transition and its implications for future development. Furthermore, we also made a comparison study on the population pyramid of Bangladesh over the years 1951-2023.
